\documentclass[11pt,a4paper]{article}

\usepackage[utf8]{inputenc}
\usepackage[T1]{fontenc}
\usepackage[spanish,es-noquoting,es-nolists,es-noshorthands]{babel}
\usepackage{tgpagella}
\usepackage{inconsolata}
\usepackage{microtype}
\usepackage[a4paper,top=2.2cm,bottom=2.2cm,left=2.4cm,right=2.4cm]{geometry}
\usepackage{xcolor}
\usepackage{titlesec}
\usepackage{enumitem}
\usepackage{booktabs}
\usepackage{tabularx}
\usepackage{array}
\usepackage{needspace}
\usepackage{lastpage}
\usepackage{fancyhdr}
\setlength{\headheight}{14pt}
\usepackage[most]{tcolorbox}
\usepackage[colorlinks=true,linkcolor=arcablue,urlcolor=arcablue,citecolor=arcablue]{hyperref}

% ---------------------------------------------------------------- estilo
\definecolor{arcablue}{HTML}{1F3A63}
\definecolor{arcagrey}{HTML}{5A6472}
\definecolor{arcalight}{HTML}{EEF2F7}
\definecolor{arcarule}{HTML}{C9D3E0}

\titleformat{\section}{\normalfont\large\bfseries\color{arcablue}}{\thesection.}{0.6em}{}
\titleformat{\subsection}{\normalfont\normalsize\bfseries}{\thesubsection}{0.6em}{}
\titlespacing{\section}{0pt}{1.3ex plus .4ex minus .2ex}{0.8ex}
\titlespacing{\subsection}{0pt}{1.2ex plus .3ex}{0.6ex}

\setlist[itemize]{leftmargin=1.2em,itemsep=0.15ex,topsep=0.4ex,parsep=0pt}
\setlist[enumerate]{leftmargin=1.5em,itemsep=0.25ex,topsep=0.4ex,parsep=0pt}

\pagestyle{fancy}
\fancyhf{}
\renewcommand{\headrulewidth}{0.4pt}
\renewcommand{\headrule}{\hbox to\headwidth{\color{arcarule}\leaders\hrule height \headrulewidth\hfill}}
\fancyhead[L]{\footnotesize\color{arcagrey}Modelos Generativos Profundos · MIA · ICAI}
% \fancyhead[R]{\footnotesize\color{arcagrey}Práctica: ARCA}
\fancyfoot[C]{\footnotesize\color{arcagrey}\thepage\ de \pageref{LastPage}}

\newcommand{\code}[1]{{\small\texttt{#1}}}
\newcommand{\file}[1]{{\small\texttt{\textcolor{arcablue}{#1}}}}

\newtcolorbox{minimos}[1]{
  enhanced, breakable, colback=arcalight, colframe=arcablue,
  boxrule=0pt, leftrule=2.5pt, arc=1pt, left=8pt, right=8pt, top=6pt, bottom=6pt,
  fonttitle=\bfseries\small, title={#1}, coltitle=arcablue,
  attach boxed title to top left={xshift=8pt,yshift=-2pt},
  boxed title style={colback=arcalight,boxrule=0pt,arc=0pt}
}

\newtcolorbox{nota}{
  enhanced, breakable, colback=white, colframe=arcarule,
  boxrule=0.6pt, arc=1pt, left=8pt, right=8pt, top=5pt, bottom=5pt
}

\setlength{\parindent}{0pt}
\setlength{\parskip}{0.55ex}

% ---------------------------------------------------------------- documento
\begin{document}

\begin{center}
  {\Large\bfseries\color{arcablue} Instrucciones de funcionamiento de la DGX}\\[0.6ex]
  {\large Entrenamiento e inferencia}\\[1.2ex]
  {\small Modelos Generativos Profundos · Máster en Inteligencia Artificial · ICAI · Curso 2026-2027}\\[0.8ex]
\end{center}


Para la realización de las prácticas y proyectos de este curso, se han modificado el tipo 
de sesiones disponibles en la DGX, adecuándose a las necesidades que podáis tener tanto para el entrenamiento
como para la inferencia.

\section{Recursos disponibles}

La DGX dispone de 8 GPUs Nvidia H200 de 141GB de VRAM. Se usa el sistema de GPU multi-instancia (MIG),
cuyas nuevas instancias disponibles son las siguientes:

\begin{itemize}
  \item \textbf{46 instancias de 1/7 de GPU (16GB)}. Están pensadas para entrenamiento e inferencia
  de modelos pequeños, especialmente para el desarrollo del agente ARCA.
  \item \textbf{1 instancia de 1/2 de GPU (71GB)}. Pensada para aquellos equipos que quieran realizar 
  \textit{fine-tuning} e inferencia de modelos con mayor capacidad, principalmente para el desarrollo
  del proyecto final. \textbf{Tiene un máximo de 7 días por sesión}. Cuando se lanza un trabajo en esta isntancia,
  queda ocupada y los trabajos que la requieran quedarán en cola de espera.
  \item \textbf{1 instancia de 1/2 de GPU (71GB)}. Controlada por los profesores solamente y sin límite
  de tiempo. Está pensada para proporcionar un entorno de inferencia común que pueda servir a todos los grupos
  para:
  \begin{itemize}
    \item Diseño de aplicaciones de agentes que no requieran un entrenamiento adicional
    \item Generación de trazas de razonamiento para el arranque en frío del SFT para el agente ARCA.
  \end{itemize}
\end{itemize}

\section{Servicio de inferencia común}

Una instancia gestionada por los profesores está activa de manera indefinida para todos los grupos 
de la asignatura. 

En esta instancia se ha desplegado un servidor de inferencia vLLM con el modelo \href{https://huggingface.co/Qwen/Qwen3.8-27B-FP8}{Qwen3.8-27B-FP8},
con capacidades agénticas y de razonamiento bastante avanzadas. El servidor optimiza el servicio 
de inferencia para servir a múltiples usuarios de manera simultánea con un rendimiento alto (compilación,
batching, KV cache, kernel fusion...). Estos son los ajustes y prestaciones más interesante de cara al usuario:

\begin{itemize}
\item Longitud de contexto máxima de 65.536 tokens, incluyendo el prompt.
\item Máximo de 5 secuencias simultáneas en proceso.
\item Alrededor de 50 tokens por segundo, dependiendo de la longitud de contexto y número de requests simultáneas.
\item Entre 1 y 7s de tiempo hasta el primer token, dependiendo de la longitud de contexto y número de requests simultáneas
\end{itemize}

Para acceder a este servicio, basta con apuntar a la URL \url{http://127.0.0.1:8000} siempre \textbf{dentro de la DGX}.
El servidor no es accesible desde redes externas. No todos los endpoints están expuestos, estando 
disponibles:

\begin{itemize}
  \item \code{GET:/health}: check de estado del servidor.
  \item \code{GET:/v1/models} modelos disponibles.
  \item \code{POST:/v1/chat/completions}: generación de texto en formato chat.
  \item \code{POST:/v1/completions}: generación de texto.
  \item \code{GET:/metrics}: check de métricas del servidor
\end{itemize}

\pagebreak

\subsection{Ejemplo de uso}

A continuación se muestra un ejemplo de uso del endpoint \code{POST:/v1/chat/completions} con \texttt{curl}:

\begin{verbatim}
 curl -X POST http://127.0.0.1:8000/v1/chat/completions \
  -H "Content-Type: application/json" \
  -d '{
    "model": "qwen3.8-27b-fp8",
    "messages": [
      {
        "role": "user",
        "content": "Reply with the word READY."
      }
    ],
    "max_tokens": 64,
    "chat_template_kwargs": {
      "enable_thinking": false
    }
  }' 
\end{verbatim}

Nótese el nombre del modelo \texttt{qwen3.8-27b-fp8} y el parámetro \texttt{chat\_template\_kwargs},
que permite desactivar la generación de trazas de razonamiento.

\textbf{Los ajustes de muestreo deberían ser los de por defecto usando el modo con razonamiento}. Para usar
el modo sin razonamiento, consultar los parámetros en \href{https://huggingface.co/Qwen/Qwen3.8-27B-FP8}{la ficha del modelo}.

Los frameworks de aplicaciones de agentes tienen normalmente soporte para estas APIs tipo OpenAI. Por ejemplo,
LangChain permite establecer conexión mediante un cliente definido como sigue:

\begin{verbatim}
  ChatOpenAI(
        model="qwen3.8-27b-fp8",
        base_url="http://localhost:8000/v1",
        api_key="EMPTY",
    )
\end{verbatim}

donde \code{api\_key} puede ser cualquier valor, ya que el servidor no requiere autenticación.

\end{document}